\documentclass[10pt]{article}
%% korr-preamble.tex -- gemeinsame Präambel für die Korrekturlisten LNM 1184
%% Makros wie im Buch (LN-Book-Springer/preamble), damit die Formeln
%% genauso aussehen wie im Proof.
\usepackage[T1]{fontenc}
\usepackage[utf8]{inputenc}
\usepackage{newtxtext}
\usepackage[varvw]{newtxmath}
\usepackage[scaled=.9]{helvet}
\usepackage[varqu,scaled=.95]{zi4}
\usepackage[ngerman]{babel}
\usepackage{csquotes}
\usepackage[a4paper,landscape,margin=1.6cm,top=1.8cm,bottom=1.8cm]{geometry}
\usepackage{microtype}
\usepackage{array,longtable,booktabs}
\usepackage{ragged2e}
\usepackage[table]{xcolor}
\usepackage{xspace}
\usepackage{enumitem}
\setlist{nosep,leftmargin=*}
\usepackage[nointegrals]{wasysym}
\usepackage{newunicodechar}
\usepackage{nicefrac}
\usepackage{mathtools}
\usepackage{mleftright}
\mleftright
\usepackage[hidelinks]{hyperref}
\setlength{\parindent}{0pt}
\setlength{\parskip}{.5em}
\setlength{\emergencystretch}{3em}
\renewcommand{\arraystretch}{1.25}
\definecolor{codegray}{gray}{0.35}

%% ---- Makros aus dem Buch -------------------------------------------
\DeclarePairedDelimiterX{\norm}[1]{\lVert}{\rVert}{\ifblank{#1}{\:\cdot\:}{#1}}
\DeclarePairedDelimiterX{\abs}[1]{\lvert}{\rvert}{\ifblank{#1}{\:\cdot\:}{#1}}
\DeclarePairedDelimiterX{\dualp}[1]{\langle}{\rangle}{\ifblank{#1}{\,\cdot\,,\cdot\,}{#1}}
\newcommand{\interval}[1]{\left[ #1 \right]}
\newcommand{\ointerval}[1]{\left] #1 \right[}
\newcommand{\rointerval}[1]{\left[ #1 \right[}
\newcommand{\lointerval}[1]{\left] #1 \right]}
\newcommand{\tint}{\textstyle\int}
\AtBeginDocument{%
\let\ORGvarepsilon=\varepsilon
\let\epsilon=\epsilon
\let\epsilon=\ORGvarepsilon
\let\ORGvarrho=\varrho
\let\rho=\rho
\let\rho=\ORGvarrho
\let\ORGvartheta=\vartheta
\let\theta=\theta
\let\theta=\ORGvartheta
\let\ORGvarphi=\varphi
\let\phi=\phi
\let\phi=\ORGvarphi
\let\ORGvarleq=\leqslant
\let\leqslant=\leq
\let\leq=\ORGvarleq
\let\ORGvargeq=\geqslant
\let\geqslant=\geq
\let\geq=\ORGvargeq
\let\ORGvarpreccurlyeq=\preccurlyeq
\let\preccurlyeq=\preceq
\let\preceq=\ORGvarpreccurlyeq
\let\setminus\smallsetminus			% A - B fuer Mengen
\renewcommand{\P}{\mathfrak{P}}		% Achtung: \P hat andere Bedeutung in LaTeX
}
\newcommand{\N}{\mathbb{N}}		% Natuerliche Zahlen
\newcommand{\Z}{\mathbb{Z}}		% Ganze Zahlen
\newcommand{\Q}{\mathbb{Q}}		% Rationale Zahlen
\newcommand{\R}{\mathbb{R}}		% Reelle Zahlen
\newcommand{\C}{\mathbb{C}}		% Komplexe Zahlen
\newcommand{\K}{\mathbb{K}}		% Koerperzeichen
\newcommand{\T}{\mathbb{T}}		% Torus
\DeclareMathOperator{\Id}{Id}
\DeclareMathOperator{\sign}{sign}				% signum-Operator
\newcommand{\im}{\ensuremath{\mathrm{i}}}		% imaginäre Konstante i
\newcommand{\iu}{\ensuremath{\mathrm{i}}}
\newcommand{\eu}{\ensuremath{\mathrm{e}}}
\renewcommand*{\Re}{\mathop{}\!\mathrm{Re}}		% Realteil
\renewcommand*{\Im}{\mathop{}\!\mathrm{Im}}		% Imaginärteil
\newcommand*{\supp}{\mathop{}\!\mathrm{supp}\,}  % Träger
\newcommand*{\diff}[1]{\mathop{}\!\mathrm{d}{#1}}	% \diff{s} = ds, 
\newcommand*{\ds}{\mathop{}\!\mathrm{d}{s}}       	% \ds = ds, 
\newcommand*{\dg}{\mathop{}\!\mathrm{d}{g}}       	% \dg = dg, 
\newcommand*{\dt}{\mathop{}\!\mathrm{d}{t}}       	% \dt = dt, 
\newcommand*{\dx}{\mathop{}\!\mathrm{d}{x}}       	% \dx = dx, 
\newcommand*{\dy}{\mathop{}\!\mathrm{d}{y}}       	% \dy = dy, 
\newcommand*{\dr}{\mathop{}\!\mathrm{d}{r}}       	% \dr = dr, 
\newcommand*{\dm}{\mathop{}\!\mathrm{d}{m}}       	% \dm = dm, 
\newcommand*{\du}{\mathop{}\!\mathrm{d}{u}}       	% \du = du, 
\DeclareMathOperator{\esssup}{ess\,sup}
\newcommand{\LE}{\mathcal{L}(E)} 		% 
\newcommand{\BH}{\mathcal{B}(H)} 		% 
\newcommand{\BA}{\mathfrak{A}}
\newcommand{\ZE}{\mathcal{Z}(E)} 	% \ZE  Zentrum von E
\newcommand{\LH}{\mathcal{L}(H)} 
\let\originalL\L
\AtBeginDocument{%
  \DeclareRobustCommand{\L}[1]{\mathcal{L}(#1)}%
}
\newcommand{\Fix}[1]{\mathop{}\mathrm{Fix}\left(#1\right)}		% \Fix{T} Fixraum
\newcommand{\rank}{\mathop{}\!\mathrm{rank}\,}		% 
\newcommand{\Kern}[1]{\mathop{}\mathrm{ker}(#1)}
\newcommand{\Image}[1]{\mathop{}\mathrm{im}(#1)}
\newcommand{\image}{\mathop{}\!\mathrm{im}\,}
\renewcommand{\dim}{\mathop{}\!\mathrm{dim}\,}
\newcommand*{\eg}{e.g.,\xspace}
\newcommand*{\ie}{i.e.,\xspace}
\newcommand*{\resp}{resp.\xspace}
\newcommand*{\vs}{vs.\xspace}
\newcommand*{\cf}{cf.\xspace}
\newcommand*{\etc}{etc.\xspace}
\newcommand{\CA}{$\mathrm{C}^{*}$}	% C*-Algebra
\newcommand{\WA}{$\mathrm{W}^{*}$}	% W*-Algbera
\newcommand{\AW}{$\mathrm{AW}^{*}$}	% AW*-Algbera
\newcommand*{\TT}{\mathcal{T}}		% C_0-Semigroup
\newcommand*{\RR}{\mathcal{R}}		% Pseudo Resolvente
\newcommand*{\SG}{\mathcal{S}}	    % Semigroup
\newcommand*{\UG}{\mathcal{U}}	    % Implemented Semigroup
\newcommand*{\F}{\mathcal{F}}		% F-product
\newcommand{\mfrac}[2]{%
    {\raise.5ex\hbox{$\scriptstyle #1$}}%
    \kern.025em/\kern.025em%
    {\lower.25ex\hbox{$\scriptstyle #2$}}}
\newcommand*{\sfrac}[2]{\leavevmode\kern.1em%
       \raise.5ex\hbox{\small #1}%
       \kern.025em/\kern.025em%
       \lower.25ex\hbox{\small #2}}
\newcommand{\1}{\mathbf{1}}
\newcommand{\HO}{\mathcal{H}(O)} 		%  für Lotz
\newcommand{\WO}{\mathcal{W}(O)} 		%  für Lotz
\newcommand*{\ddp}{\mathop{}\!\mathrm{d}{p}}       	% \ddp = dp, 
\newcommand*{\dN}{\mathop{}\!\mathrm{d}{N}}       	% \dN = dN, 
\newcommand*{\ddt}{\mathop{}\!\mathrm\fraq{d}{dt}}   % \ddt = d/dt
\newcommand*{\lnm}{LNM1184: } % \LMN-BookCode
\newcommand{\closure}[1]{#1^{-}}
%% Hilfsmatrizen aus Teil I, Kap. 1
\newcommand{\ABC}{\begin{pmatrix} 0 & 1 \\ -1 & 0 \end{pmatrix}}
\newcommand{\BBC}{\begin{pmatrix} \cos(t) & \sin(t) \\ - \sin(t) & \cos(t) \end{pmatrix}}
\newcommand{\CBC}{\begin{pmatrix} 1 & 1 \\ -1 & -1 \end{pmatrix}}
\newcommand{\DBC}{\begin{pmatrix} 1+t & t \\ - t & 1-t \end{pmatrix}}
\newcommand{\EBC}{\begin{pmatrix} 0 & 1 \\ 1 &0 \end{pmatrix}}
\newcommand{\FBC}{\begin{pmatrix} \cosh(t) & \sinh(t) \\ \sinh(t) & \cosh(t) \end{pmatrix}}

%% ---- Darstellung der Listen ----------------------------------------
%% Quelltext, der sich nicht als Formel setzen lässt (Labels, \ref, Dateinamen)
\newcommand{\code}[1]{\leavevmode{\color{codegray}\ttfamily\small #1}}
\newcommand{\sichA}{\CIRCLE}      % ● hoch
\newcommand{\sichB}{\LEFTcircle}  % ◐ mittel
\newcommand{\sichC}{\Circle}      % ○ niedrig
%% Querverweise im Buch: Label-Namen statt Nummern
\newcommand{\lbl}[1]{\textsf{\footnotesize\color{codegray}\detokenize{#1}}}

\hypersetup{pdftitle={Mathematische Fehler in Beweisen: Part C (Teil III)}}
\pagestyle{plain}
\begin{document}
% Erzeugt aus math-error-part-c.md: Formeln gesetzt statt Quelltext.
% Grau/Schreibmaschine = Quelltext, der sich nicht als Formel setzen laesst (Makronamen, Labels).
{\Large\bfseries Mathematische Fehler in Beweisen: Part C (Teil III)\par}
\smallskip
\emph{One-parameter Semigroups of Positive Operators}, 2. Aufl. (Springer LNM 1184). Stand: 14.10.2026.

\textbf{Umfang:} Beweise, Rechnungen und Beispiele in Teil III, Kap. 1–4, inhaltlich geprüft im LaTeX-Quelltext (Branch \code{fahnenkorrektur}): Vorzeichen und Konstanten, Ungleichungsrichtungen, Voraussetzungen und Quantoren, zitierte Hilfsaussagen, Lücken, Beispiele. Reine Satz- und Notationsfragen stehen in \code{korrekturen-mathe-part-c.md}; wo ein Befund dort schon steht und hier ergänzt wird, ist das vermerkt („Ergänzt korrekturen-mathe“).

\textbf{Wichtig:} Die Befunde sind von Claude nachgerechnet, aber nicht von einem zweiten Mathematiker geprüft. Bitte jeden Punkt selbst verifizieren, bevor er umgesetzt wird. \textbf{Am Quelltext ist nichts geändert.}

\textbf{Spalten:} S. $=$ gedruckte Buchseite (über das nächstgelegene Label ermittelt, $\ensuremath{\pm}1$ Seite). Z. $=$ Zeile in der .tex-Datei. Sicherheit: \sichA{} sicher falsch $\ensuremath{\cdot}$ \sichB{} sehr wahrscheinlich falsch $\ensuremath{\cdot}$ \sichC{} zweifelhaft, bitte prüfen.

Insgesamt \textbf{57 Punkte}: 20 \sichA{}, 18 \sichB{}, 19 \sichC{}.

\section*{Befunde}
\subsection*{Teil III, Kap. 1 – Basic Results on Banach Lattices and Positive Operators (\code{part-c/chap-c1.tex})}
\begin{longtable}{@{}rrc>{\RaggedRight\arraybackslash}p{3.0cm}>{\RaggedRight\arraybackslash}p{11.4cm}>{\RaggedRight\arraybackslash}p{7.4cm}@{}}
\toprule
\textbf{S.} & \textbf{Z.} & \textbf{} & \textbf{Stelle} & \textbf{Befund} & \textbf{Vorschlag} \\
\midrule
\endfirsthead
\toprule
\textbf{S.} & \textbf{Z.} & \textbf{} & \textbf{Stelle} & \textbf{Befund} & \textbf{Vorschlag} \\
\midrule
\endhead
\bottomrule
\endlastfoot
275 & 529 & \sichB{} & Beweis der Proposition in 1.8 (Signum-Operator), Fortsetzung auf das Band & Ergänzt korrekturen-mathe. In $S_{f}h = \sup \{ S_{f}g \colon g \in [0,h] \cap E_{\vert f\vert} \}$ wird $S_f$ durch sich selbst definiert (zirkulär). Mit $S_0$ ist die Formel korrekt: $S_0$ ist auf dem Ideal $E_{\vert f\vert}$ definiert, wegen (ii) gilt $S_{0}g \leq g \leq h$ für positives $S_0$, also existiert das Supremum in dem ordnungsvollständigen E, und weil $E_{\vert f\vert}$ ordnungsdicht im Band $\{f\}^{dd}$ ist, entsteht eine positive Fortsetzung. Dieselbe Formel wird danach auf die positiven Teile $S_1$, \dots{}, $S_4$ angewendet (Z. 533), sie ist also für ein beliebiges positives $S_0$ gemeint. & $S_{f}h = \sup \{ S_{0}g \colon g \in [0,h] \cap E_{\vert f\vert} \}$ für $h \in \{f\}^{dd}_{+}$ \\[2pt]
276 & 554 & \sichB{} & Proposition in 1.9 (Zentrum), Bedingung (b) & Die Proposition ist ausdrücklich für (komplexe) Banachverbände formuliert, aber $\pm T \leq \|T\| \Id$ ergibt nur für reelle T einen Sinn, weil die Ordnung auf $\mathcal{L}(E)$ nach Z. 476 nur für reelle Operatoren erklärt ist. Gegenbeispiel: $T = \iu\,\Id$ erfüllt (a) und (c), aber nicht (b). Damit ist (b) im komplexen Fall nicht zu (a), (c) äquivalent. Ebenso enthält das Zentrum $\mathcal{Z}(E) \cong C(K)$ im komplexen Fall nichtreelle Operatoren. & Für den komplexen Fall (b) ersetzen durch $\vert Tf\vert \leq \|T\|\,\vert f\vert$ für alle $f \in E$. Das ist im reellen Fall äquivalent zu $\pm T \leq \|T\| \Id$. Alternativ (b) ausdrücklich auf reelle T beschränken. \\
\end{longtable}
\subsection*{Teil III, Kap. 2 – Characterization of Positive Semigroups on Banach Lattices (\code{part-c/chap-c2.tex})}
\begin{longtable}{@{}rrc>{\RaggedRight\arraybackslash}p{3.0cm}>{\RaggedRight\arraybackslash}p{11.4cm}>{\RaggedRight\arraybackslash}p{7.4cm}@{}}
\toprule
\textbf{S.} & \textbf{Z.} & \textbf{} & \textbf{Stelle} & \textbf{Befund} & \textbf{Vorschlag} \\
\midrule
\endfirsthead
\toprule
\textbf{S.} & \textbf{Z.} & \textbf{} & \textbf{Stelle} & \textbf{Befund} & \textbf{Vorschlag} \\
\midrule
\endhead
\bottomrule
\endlastfoot
281 & 177 & \sichC{} & Example 2.1.5 (i) & Die Behauptung „there exists $x \in (0,1)$ such that $f(x) = \sup f = \|f^{+}\|_{\infty}$“ ist für $f \leq 0$ falsch: Dann ist das Supremum 0 und wird u. U. nur in den Randpunkten angenommen (z. B. $f(x) = -\sin(\pi x) \in D(A)$). Für diesen Fall braucht man $\phi = 0 \in \diff{N^{+}}(f)$ (wie in (iii), Z. 250). & Fallunterscheidung ergänzen: „If $f^{+} = 0$, then $0 \in \diff{N^{+}}(f)$. Otherwise there exists \dots{}“ \\[2pt]

281 & 187 & \sichA{} & Example 2.1.5 (i), Definition von $f_{0}$ & Mit den festen Integrationsgrenzen $\int_{0}^{1}$ ist $f_{0}$ eine Linearkombination von $\eu^{x}$ und $\eu^{-x}$, also $f_{0} - f_{0}'' = 0$ statt $= g$. Auch mit $\int_{0}^{x}$ in beiden Integralen ergäbe die Formel $f_{0} - f_{0}'' = -g$. & $f_{0}(x) = \frac{1}{2} \left[\eu^{x} \int_{x}^{1} \eu^{-y}g(y) \dy + \eu^{-x} \int_{0}^{x} \eu^{y}g(y) \dy \right]$ ($= \frac12\int_0^1 \eu^{-\vert x-y\vert}g(y)\dy$); damit gilt $f_{0} - f_{0}'' = g$. \\[2pt]
284 & 358 & \sichB{} & Beweis von Example 2.1.5 (v) & „$C_{c}^{\infty}(\Omega)$ \dots{} is dense in $E \cap C_{c}^{\infty}(\Omega)$“ ist eine Leerformel; gebraucht wird die Dichtheit in $E$. Außerdem wird $g$ in Z. 359 ohne Einführung benutzt. Erst damit liefert $C_{c}^{\infty}(\Omega) \subseteq (\Id - A)D(A)$ die Dichtheit von $(\Id - A)D(A)$. & „\dots{} is dense in $E$. Let $g \in C_{c}^{\infty}(\Omega)$. We show that there exists \dots{}“ \\[2pt]
286 & 547 & \sichB{} & Beweis von Lemma 2.1.10 & Aus $\psi_{0}(f) = \sup \phi([0,f] \cap E_{(g^{-})})$ folgt nur $\langle g^{-},\psi \rangle \geq \langle g^{-},\phi \rangle$, nicht Gleichheit. Außerdem liefert Z. 537 nur $\langle g^{-},\phi \rangle \neq 0$; für $> 0$ muss man ggf. $\phi$ durch $-\phi$ ersetzen. Analog ist in Z. 551 $\langle u,\psi \rangle = 0$ zu begründen ($[0,u] \subseteq E_{u} \subseteq \ker\phi$), nicht über $\leq \langle u,\phi \rangle$. & In Z. 537 „(replacing $\phi$ by $-\phi$ if necessary) $\langle g^{-},\phi \rangle > 0$“; Z. 547 $\langle g^{-},\psi \rangle \geq \langle g^{-},\phi \rangle > 0$; Z. 551 $\langle u,\psi \rangle = 0$. \\[2pt]
287 & 606 & \sichC{} & Beweis von Theorem 2.1.11, $(a) \ensuremath{\Rightarrow} (d)$ & Gezeigt wird: Ist $(A+\lambda)f \not\geq 0$ für ein $f \geq 0$, so ist $\lambda \leq \|A\|$. Das ergibt $A + \lambda \geq 0$ nur für $\lambda > \|A\|$, für $\lambda = \|A\|$ folgt kein Widerspruch. Es fehlt der Grenzübergang $\lambda \downarrow \|A\|$ (Abgeschlossenheit von $E_{+}$). & Am Ende ergänzen: „Hence $A + \lambda \geq 0$ for all $\lambda > \|A\|$, and letting $\lambda \downarrow \|A\|$ gives (d) since $E_{+}$ is closed.“ \\[2pt]
306 & 1873 & \sichA{} & Beweis von Proposition 2.4.8, $(b) \ensuremath{\Rightarrow} (a)$ & Mit $f \coloneq R(\lambda,A)g \in D(A)_{+}$ und $\phi \coloneq R(\lambda,B)'\psi \in D(B')_{+}$ (Z. 1876) entsteht $\langle f,B'\phi\rangle - \langle Af,\phi\rangle$. Bedingung (b) gilt aber für $f \in D(B)_{+}$, $\phi \in D(A')_{+}$ und liefert diese Ungleichung nicht; die Rollen von A und B sind vertauscht. & $f \coloneq R(\lambda,B)g \in D(B)_{+}$, $\phi \coloneq R(\lambda,A)'\psi \in D(A')_{+}$; dann $\langle R(\lambda,B)g - R(\lambda,A)g,\psi\rangle = \langle f,(\lambda-A')\phi\rangle - \langle(\lambda-B)f,\phi\rangle = \langle Bf,\phi\rangle - \langle f,A'\phi\rangle \leq 0$. \\[2pt]
309 & 2071 & \sichB{} & Beweis von Example 2.4.14 & Ergänzt korrekturen-mathe. Die Annahme „$\alpha = \Re(\beta)$ und $\mu = \vert\nu\vert$“ ist zu speziell: In Z. 2119 wird die Hinlänglichkeit von $\Re(\beta) \leq \alpha$, $\vert\nu\vert \leq \mu$ allgemein beansprucht. Die Rechnung in Z. 2109–2110 benutzt ohnehin nur diese Ungleichungen. & „We first assume that $\Re(\beta) \leq \alpha$ and $\vert\nu\vert \leq \mu$.“ \\[2pt]
312 & 2218 & \sichC{} & Beweis von Lemma 2.4.18 & Für $\omega_{0}(A) < -1$ dürfen $\lambda, \mu \geq \omega_{0}(A)+1$ auch $\leq 0$ sein. Dann ist $\|R(\lambda,A)\| \leq M/\lambda$ sinnlos, und $Af \leq \mu^{2}R(\mu,A)f - \mu f \leq \mu^{2}R(\mu,A)f$ (Z. 2224) braucht $\mu \geq 0$. & $\lambda, \mu \geq \max\{\omega_{0}(A)+1, 1\}$ (ebenso Z. 2219, 2231). \\[2pt]
312 & 2278 & \sichA{} & Beweis von Theorem 2.4.17 & $\sup_{\alpha}\Re(\eu^{\im \alpha}B) = \vert B\vert = \vert B_{0}\vert + \vert M\vert$, nicht $\vert B_{0}\vert$. Nach Anwendung der Bandprojektion auf $\mathcal{Z}(E)^{d}$ steht links $\Re(\eu^{\im\alpha}B_{0})$. & $\vert B_{0}\vert = \sup_{\alpha \in \R}\Re(\eu^{\im \alpha}B_{0}) \leq C_{0}$ \\[2pt]
312 & 2284 & \sichA{} & Beweis von Theorem 2.4.17 (Schluss) & Gezeigt wird $\Re((\sign \bar{f})Bf) \leq C\vert f\vert$, also dass $(\eu^{tC})$ die Halbgruppe $(\eu^{tB})$ dominiert. Das ist die Voraussetzung. Für die Minimalität muss aber $(\eu^{tC})$ die Halbgruppe $(\eu^{tA})$ mit $A = \vert B_{0}\vert + \Re(M)$ dominieren. & $\Re((\sign \bar{f})Af) = \Re((\sign \bar{f})\vert B_{0}\vert f) + \Re(M)\vert f\vert \leq C_{0}\vert f\vert + N\vert f\vert = C\vert f\vert$; Z. 2286: „\dots{} that $(\eu^{tA})_{t \geq 0}$ is dominated by $(\eu^{tC})_{t \geq 0}$.“ (Alternativ: $A \leq C$, dann Proposition 2.4.8.) \\[2pt]
317 & 2551 & \sichC{} & Remark 2.5.7 (ii) & Für reellwertige konvexe Funktionen $\Theta \colon E \to \R_{+}$ ist die einseitige Gateaux-Differenzierbarkeit elementar; die Voraussetzung „order continuous norm“ spielt dann keine Rolle. Gemeint ist offenbar (wie im Beweis von Proposition 2.5.6) eine konvexe Abbildung mit Werten im Banachverband. & bitte prüfen, z. B. „every convex function $\Theta$ from a vector space into $E$“. \\[2pt]
318 & 2603 & \sichC{} & Text vor Corollary 2.5.8 & „Kato's equality ({\ie{}} the equality holds in (2.5.13))“: Gleichheit in (2.5.13) bedeutet, dass T(t) ein Verbandshomomorphismus ist. Die Aussage wäre damit tautologisch. Gemeint ist Gleichheit in Katos Ungleichung für den Generator (vgl. (2.2.10), Corollary 2.5.8 (c)). & {(\ie the equality holds in Eq.~(\lbl{eq:c2-2.10}))} \\[2pt]
320 & 2800 & \sichA{} & Beweis von Theorem 2.5.13, $(e) \ensuremath{\Rightarrow} (d)$ & „$\vert g\vert \leq \vert f\vert$ $\ensuremath{\Rightarrow}$ there exists a band projection P such that $Pf = g$“ ist falsch, z. B. für $g = f/2$, $f \neq 0$ (aus $Pf = f/2$ folgt $Pf = P^{2}f = Pf/2$, also $f = 0$). Richtig ist $g = Sf$ mit $S \in \mathcal{Z}(E)$, $\vert S\vert \leq \Id$. & bitte prüfen; z. B.: S durch Linearkombinationen von Bandprojektionen approximieren (Freudenthal); mit (e) und der Abgeschlossenheit von A folgt $g = Sf \in D(A)$. \\[2pt]
322 & 2898 & \sichA{} & Beweis von Lemma 2.5.17 & Ergänzt korrekturen-mathe. Für das positive Minimumprinzip (P) muss $\phi$ positiv sein: $\phi \in E'_{+}$, nicht nur $\phi \in E'$. Außerdem folgt $\langle Af,\phi\rangle = 0$ aus $Af \in \{f\}^{dd}$ nur für positives $\phi$, denn dann verschwindet $\phi$ auf $\{f\}^{dd} = \overline{E_{f}}$. & $0 \leq f \in D$, $\phi \in E'_{+}$ such that $\langle f,\phi \rangle = 0$. \\
\end{longtable}
\subsection*{Teil III, Kap. 3 – Spectral Theory of Positive Semigroups on Banach Lattices (\code{part-c/chap-c3.tex})}
\begin{longtable}{@{}rrc>{\RaggedRight\arraybackslash}p{3.0cm}>{\RaggedRight\arraybackslash}p{11.4cm}>{\RaggedRight\arraybackslash}p{7.4cm}@{}}
\toprule
\textbf{S.} & \textbf{Z.} & \textbf{} & \textbf{Stelle} & \textbf{Befund} & \textbf{Vorschlag} \\
\midrule
\endfirsthead
\toprule
\textbf{S.} & \textbf{Z.} & \textbf{} & \textbf{Stelle} & \textbf{Befund} & \textbf{Vorschlag} \\
\midrule
\endhead
\bottomrule
\endlastfoot
327 & 49 & \sichA{} & Theorem 3.1.1 (iii) & Die Aussage $h_{0} \in E_{\vert h\vert}$ (Hauptideal) ist zu stark. Gegenbeispiel: Rotationshalbgruppe mit Periode 1 auf $L^1(\ensuremath{\Gamma})$. Hier ist $T(1) = \Id $, ker A besteht nur aus Konstanten, und ein unbeschränktes $h_0 \ensuremath{\geq} 0$ aus $L^1$ liegt in keinem $E_{\vert h\vert}$. Mit h $= \ensuremath{\int}_0^1 T(s)h_0$ ds (wie in Part-II, Proposition 3.1.5) folgt nur, dass $h_0$ im abgeschlossenen Ideal liegt, das h erzeugt. Das passt auch zur Nullideal-Formulierung in (iv) und zur Trägerinklusion in Part II. & $h_{0} \in \overline{E_{\vert h\vert}}$ (abgeschlossenes Ideal, das von h erzeugt wird) \\[2pt]
328 & 92 & \sichA{} & Beweis von Theorem 3.1.2, Eq. (3.1.3) & Nach Eq. (3.1.2) gilt $R(\lambda_{0},A)^{n+1}f = \frac{1}{n!}\int s^n \eu^{-\lambda_0 s}T(s)f\ds$, deshalb fehlt in Z. 92 und Z. 93 der Faktor $s^n$. Erst mit ihm ergibt die Summation in Z. 94 $\exp((\lambda_{0}-\mu)s)$. & Z. 92: $\int_{0}^{\infty} \frac{1}{n!}(\lambda_{0}-\mu)^n s^n\exp(-\lambda_{0}s)\langle T(s)f,\phi\rangle\ds$; Z. 93: $\sum_{n=0}^{\infty} \frac{1}{n!}(\lambda_{0}-\mu)^n s^n\exp(-\lambda_{0}s)\langle T(s)f,\phi\rangle$ \\[2pt]
%333 & 374 & \sichC{} & Example 3.2.5 (i) & In „not the resolvent of an operator on $G$“ ist der Raum falsch: $R_F$ wirkt auf $G_F$, und gezeigt wird wegen $0 \in P\sigma(R_{\F}(\lambda))$, dass $R_F$ keine Resolvente eines Operators auf $G_F$ ist. & {It follows that $R_{\F}$ is not the resolvent of an operator on $G_{\F}$.} \\[2pt]
334 & 445 & \sichA{} & Beweis von Proposition 3.2.6 (i) & Der Faktor $(\mu-\alpha)$ links ist falsch. Nach Eq. (3.2.8) gilt $(\lambda-\alpha)(R(\lambda)x - R(\mu)x) = (\lambda-\alpha)(\mu-\lambda)R(\mu)R(\lambda)x$, links muss also $(\lambda-\alpha)$ stehen. Wegen $(\lambda-\alpha)R(\lambda)x = x$ ist das gleich $(\mu-\lambda)R(\mu)x$, und daraus folgt $(\mu-\alpha)R(\mu)x = x$. & $(\mu-\lambda)R(\mu)x = (\lambda-\alpha)(\mu-\lambda)R(\mu)R(\lambda)x = (\lambda-\alpha)(R(\lambda)x - R(\mu)x) = x - (\lambda-\alpha)R(\mu)x$ \\[2pt]
336 & 551 & \sichC{} & Remarks nach Definition 3.2.8 (i) & Die Menge $\{\frac{1}{t}\int_0^t \ldots \colon 0 \leq t \leq t_{0}\}$ enthält $t = 0$, und dort ist $\frac{1}{t}$ nicht definiert. & $0 < t \leq t_{0}$ \\[2pt]
338 & 651 & \sichB{} & Beweis von Theorem 3.2.10 & Gewählt wird nur $\phi \in E_{1}'$ mit $\langle h_{1},\phi\rangle \neq 0$, aber nicht $\ensuremath{\varphi} \ensuremath{\geq} 0$. Später wird $\ensuremath{\varphi}_1 \ensuremath{\geq} 0$ gebraucht (Abschätzung Z. 678, Idealeigenschaft von I in Z. 680, Eq. (3.2.19)). Außerdem liefert die Kette in Z. 670 die untere Schranke nur für $\tau R_{1}(\tau)'\vert\phi\vert$. & {Next, we choose any $\phi \in (E_{1}')_{+}$ such that $\langle \vert h_{1}\vert,\phi \rangle > 0$.} (Dann ist $\ensuremath{\varphi}_1 \ensuremath{\geq} 0$ als schwach*-Limes positiver Funktionale.) Alternativ $\ensuremath{\varphi}$ in Z. 652/653 durch $\vert \ensuremath{\varphi}\vert$ ersetzen. \\[2pt]
338 & 670 & \sichA{} & Beweis von Theorem 3.2.10 & In der Ungleichungskette fehlt der Faktor $\ensuremath{\tau}$: Eq. (3.2.17) liefert $\vert h_{1}\vert \leq \tau R_{1}(\tau)\vert h_{1}\vert$, nicht $R_{1}(\tau)\vert h_{1}\vert$. Außerdem muss im letzten Term $\vert\phi\vert$ statt $\phi$ stehen. & $0 < \vert\langle h_{1},\phi \rangle\vert \leq \langle \vert h_{1}\vert,\vert\phi\vert \rangle \leq \langle \tau R_{1}(\tau)\vert h_{1}\vert,\vert\phi\vert \rangle = \langle \vert h_{1}\vert,\tau R_{1}(\tau)'\vert\phi\vert \rangle$ (mit $\ensuremath{\varphi} \ensuremath{\geq} 0$ wie im Befund zu Z. 651 entfällt $\vert \ensuremath{\varphi}\vert$) \\[2pt]
348 & 1303 & \sichA{} & Beweis von Corollary 3.3.9 & $\alpha \coloneq \frac{2\pi}{n\tau}$ ist falsch. Für eine $\ensuremath{\tau}$-periodische Halbgruppe gilt $\sigma(A) \subseteq \frac{2\pi\im}{\tau}\Z$, und dessen Untergruppen sind $\im\frac{2\pi n}{\tau}\Z$. Für $n > 1$ läge $\im\frac{2\pi}{n\tau}\Z$ nicht in $\frac{2\pi\im}{\tau}\Z$. & $\alpha \coloneq \frac{2\pi n}{\tau}$ \\[2pt]
348 & 1305 & \sichC{} & Beweis von Corollary 3.3.9 & Hier wird nur $\phi \in \ker A'$ verlangt. Für den AL-Raum $(E,\ensuremath{\varphi})$ und für die strikte Positivität nach Proposition 3.3.5 (ii) (Z. 1317 f.) muss aber $\ensuremath{\varphi} \ensuremath{\geq} 0$ sein. Das lässt sich erreichen, z. B. mit $\ensuremath{\varphi} = P'\ensuremath{\psi}$, $\ensuremath{\psi} \ensuremath{\geq} 0$, wobei P die Mittelwertprojektion $(1/\ensuremath{\tau})\ensuremath{\int}_0^\ensuremath{\tau}$ T(s) ds ist. & {there exists $\phi \in \ker A'$, $\phi \geq 0$, such that $\langle\vert h\vert,\phi\rangle = 1$} \\[2pt]
350 & 1402 & \sichA{} & Beweis von Theorem 3.3.12, Eq. (3.3.15) & Vorzeichenfehler: Aus Eq. (3.3.13) $S_{h}^{-1}AS_{h} = A+\im\alpha$ folgt $S_{h}^{-k}R(\lambda,A)S_{h}^{k} = R(\lambda-\im k\alpha,A)$, wie in Eq. (3.2.14). Die gedruckte Formel würde $R(\lambda+\im k\alpha,A) = R(\lambda-\im k\alpha,A)$ ergeben. Der Schluss in Z. 1407 bleibt richtig. h (Eigenvektor zu $i\ensuremath{\alpha}$ nach Theorem 3.3.8 (iii)) wird im Beweis nicht eingeführt. & $R(\lambda+\im k\alpha,A) = S_{h}^{k}\circ R(\lambda,A)\circ S_{h}^{-k}$ \\[2pt]
352 & 1545 & \sichB{} & Example 3.3.15 & Ergänzt korrekturen-mathe. Ein Teilchen mit Geschwindigkeitsbetrag $\vert v\vert$ hat bei Reflexion an beiden Enden die Periode $2/\vert v\vert$. Die zugehörigen Eigenwerte sind $\im\pi k\vert v\vert$ ($k \ensuremath{\in}$ Z), also $\sigma(A) = \{\im\gamma \colon n\pi v_{1} \leq \vert\gamma\vert \leq n\pi v_{2}\}$. Das bestätigt $\ensuremath{\gamma}_i = \ensuremath{\pi} v_i$. & $\gamma_{1} \coloneq \pi v_{1}$, $\gamma_{2} \coloneq \pi v_{2}$ \\[2pt]
352 & 1593 & \sichC{} & Beweis von Corollary 3.3.17 & Ohne die Normierung $s(A) = 0$ (die hier fehlt) ist $\{\lambda \in \sigma(A) \colon \Re(\lambda) \geq -1\}$ die falsche Menge: Für $s(A) < -1$ ist sie leer und gibt keine Lücke unterhalb von s(A) her. Gebraucht wird eine Halbebene links von s(A). & $\{\lambda \in \sigma(A) \colon \Re(\lambda) \geq s(A)-1\}$ \\[2pt]
354 & 1655 & \sichA{} & Beweis von Proposition 3.3.18 (ii) & m wird mit $\langle K^{m}f,\phi\rangle > 0$ gewählt. Der Integrand von $\langle R_{m}(t)f,\phi\rangle$ ist an der Stelle $t_1 =$ \dots{} $= t_m = t$ aber $\langle K^{m}S(t)f,\phi\rangle$ (wegen $T(0) = \Id $ und $S(t_m) = S(t)$). Die Irreduzibilität von K muss daher auf $S(t)f > 0$ angewandt werden; genau dafür wird in Z. 1654 $S(t)f > 0$ gezeigt. & {there exists $m \in \N$ such that $\langle K^{m}S(t)f,\phi\rangle > 0$}; in Z. 1656 $t_{1}=t_{2}=\ldots=t_{m}=t$ \\[2pt]
354 & 1680 & \sichA{} & Transportgleichung nach Proposition 3.3.18 & Vorzeichenfehler: $\lambda-(A-M) = (\Id + MR(\lambda,A))(\lambda-A)$, also $R(\lambda,A-M) = R(\lambda,A)(\Id + MR(\lambda,A))^{-1}$. Vgl. Z. 1627 mit B $= A + K$ und Minuszeichen. & $R(\lambda,A-M) = R(\lambda,A)(\Id + MR(\lambda,A))^{-1}$ \\[2pt]
354 & 1685 & \sichC{} & Transportgleichung nach Proposition 3.3.18 & Die Bedingung „provided that $s(B) > s(A)$“ passt nicht zu Z. 1681. Dort wird Proposition 3.3.18 auf $A \ensuremath{-} M$ und K angewandt, also mit der Bedingung $s(B) > s(A \ensuremath{-} M$). & {(provided that $s(B) > s(A-M)$)}; bitte prüfen \\[2pt]
356 & 1793 & \sichC{} & Beweis von Lemma 3.4.1 & Theorem 3.2.2 setzt eine stark stetige positive Halbgruppe mit Generator voraus. $(T_1(t)')$ auf $(E,\ensuremath{\psi})'$ ist im Allgemeinen nur schwach*-stetig. Die Folgerung stimmt trotzdem direkt: Die $T_1(t)'$ sind Verbandsisomorphismen, also gilt nach Eq. (3.2.5) $T_{1}(t)'\phi^{[n]} = (T_{1}(t)'\phi)^{[n]} = \eu^{\im n\alpha t}\phi^{[n]}$ und damit $A_{1}^*\phi^{[n]} = \im n\alpha\phi^{[n]}$. & Begründung direkt über die Verbandsisomorphismen statt mit Theorem 3.2.2; bitte prüfen \\[2pt]
358 & 1897 & \sichC{} & Example 3.4.4, Eq. (3.4.6) & Ergänzt korrekturen-mathe. $\frac{1}{2\pi}\int_{0}^{2\pi} f(z\cdot\eu^{\im s})\ds$ ist der Mittelwert von f über $\ensuremath{\Gamma}$ und hängt nicht von z ab. Mathematisch ist die Formel also korrekt; die Ersetzung durch $f(\eu^{\im s})$ ist gleichwertig und betrifft nur die Notation. & bitte prüfen (nur Notation, kein Rechenfehler) \\[2pt]
358 & 1902 & \sichA{} & Example 3.4.4 & $P\sigma(A) = i\R \backslash\{0\}$ ist falsch. Für einen Eigenvektor [f,g] zu $i\ensuremath{\beta}$ muss f(z $e^{it}$) $= e^{i\ensuremath{\beta}t}$ f gelten, also $\ensuremath{\beta} \ensuremath{\in}$ Z und $f = c z^\ensuremath{\beta}$. Für $\ensuremath{\beta} \ensuremath{\neq} 0$ ist der Mittelwert von f null, und g wäre Eigenvektor der Translationsgruppe auf $L^2(R)$, also $g = 0$. Für $\ensuremath{\beta} = 0$ folgt $g' = \ensuremath{-}k$ und damit $g \ensuremath{\notin} L^2$. Für $f = 0$ gibt es keine Eigenvektoren in $L^2(R)$. Ergebnis wie in Part-II, Example 3.2.13; der Schluss in Z. 1903 bleibt gültig. & $P\sigma(A) = \{\im k \colon k \in \Z,\ k \neq 0\}$ \\[2pt]
358 & 1918 & \sichB{} & Proposition 3.4.5, Eq. (3.4.7) & „for every $t \geq 0$“ ist für $t = 0$ falsch, wenn $\ensuremath{\sigma}(A) \ensuremath{\cap} R = \ensuremath{\emptyset}$. Rechts steht dann $\ensuremath{\sigma}(\Id ) \ensuremath{\cap} (0,\ensuremath{\infty}) =$ \{1\}, links $\ensuremath{\emptyset}$. Beispiel: nilpotente Linkstranslation auf $L^1[0,1]$ (Verbandshomomorphismen, $\ensuremath{\sigma}(A) = \ensuremath{\emptyset}$). Der Beweis behandelt nur $t > 0$. & $\text{for every} \quad t > 0$ \\[2pt]
360 & 2008 & \sichB{} & Kato-Gleichung Eq. (3.4.8), ebenso Z. 2011 und 2014 & Mit der Konvention $(\sign \bar{f})f = \vert f\vert$ (Section 2.5, Corollary 2.5.8 (c)) lautet Katos Gleichung $A\vert f\vert = \Re[(\sign \bar{f})Af]$. Mit $\sign f$ ist für komplexes f schon $\mu\vert f\vert = \Re[(\sign f)\mu f]$ (Z. 2011) falsch, denn das ergibt $\Re(f^2/\vert f\vert)$. Für reelles f (ab Z. 2017) stimmen beide Fassungen überein. & $A\vert f\vert = \Re [(\text{sign\,} \bar{f})Af]$; ebenso in Z. 2011 und 2014 $\text{sign } \bar{f}$ \\
\end{longtable}
\subsection*{Teil III, Kap. 4 – Asymptotics of Positive Semigroups on Banach Lattices (\code{part-c/chap-c4.tex})}
\begin{longtable}{@{}rrc>{\RaggedRight\arraybackslash}p{3.0cm}>{\RaggedRight\arraybackslash}p{11.4cm}>{\RaggedRight\arraybackslash}p{7.4cm}@{}}
\toprule
\textbf{S.} & \textbf{Z.} & \textbf{} & \textbf{Stelle} & \textbf{Befund} & \textbf{Vorschlag} \\
\midrule
\endfirsthead
\toprule
\textbf{S.} & \textbf{Z.} & \textbf{} & \textbf{Stelle} & \textbf{Befund} & \textbf{Vorschlag} \\
\midrule
\endhead
\bottomrule
\endlastfoot
369 & 68 & \sichB{} & Beweis von Theorem 4.1.1 (ii) & Aus $\rho \in R\sigma(T(\tau))$ folgt $s(A) \geq \omega_{0}(A)$ über den Spektralabbildungssatz für das Residualspektrum (es gibt $\lambda \in R\sigma(A)$ mit $\eu^{\tau\lambda} = \rho$, also $\Re\lambda = \omega_0(A)$). Das zitierte Theorem {\lbl{thm:a4-1.11}} (Datko; Charakterisierung von $\omega_0<0$ über $A\sigma(T(t_0))$) liefert diesen Schluss nicht; es wird in Z. 81 zu Recht für Datko benutzt. & {Part-I, Theorem~\lbl{thm:a3-6.3}} (Spektralabbildungssatz für P- und $R\ensuremath{\sigma}$) zitieren – bitte prüfen \\[2pt]
369 & 76 & \sichA{} & Beweis von Theorem 4.1.1 (iii) & Letzte Abschätzung $\eu^{-\alpha\tau}\int_0^{r-\tau}\eu^{-\alpha s}\langle T(s)f,\phi\rangle\ds \leq \|R(\alpha,A)f\|$ ist so falsch: korrekt ist nur $\leq \eu^{-\alpha\tau}\langle R(\alpha,A)f,\phi\rangle \leq \eu^{-\alpha\tau}\|\phi\|\,\|R(\alpha,A)f\|$ (Ersetzen von $\phi$ durch $2\phi$ erhält die Voraussetzung, aber nicht die Ungleichung; für $\alpha<0$ ist zudem $\eu^{-\alpha\tau}>1$). Für den Schluss mit Datko unerheblich. & $\leq \int_{0}^{\tau} \eu^{-\alpha s}\|T(s)f\| \ds + \eu^{-\alpha\tau}\|\phi\|\,\|R(\alpha,A)f\|$ \\[2pt]
369 & 85 & \sichC{} & Text nach Theorem 4.1.1 & „it is not known whether spectral and growth bounds of an arbitrary positive semigroup coincide“ für $L^p$, $p \neq 1,2,\infty$: Stand 1986; inzwischen von L. Weis (Proc. AMS 123, 1995) positiv beantwortet ($s(A)=\omega_0(A)$ für positive Halbgruppen auf $L^p$, $1\le p<\infty$). & ggf. „it was not known (in 1986)“ plus Hinweis auf Weis (1995) in den Notes – bitte prüfen \\[2pt]
374 & 316 & \sichA{} & Beweis von Proposition 4.1.6 & „there exists $-g \in D(A)_{+}$ such that $Ag = -f$“ ist falsch herum: aus $f = Ah$ mit $h \in D(A)_{-}$ folgt $g := -h \in D(A)_{+}$, $Ag = -f$; mit $-g \geq 0$ wäre $Ag=-f$ gerade nicht gegeben, und Z. 320 ($0 \leq T(t)g$) braucht $g \geq 0$. & {there exists $g \in D(A)_{+}$ (\ie $-g \in D(A)_{-}$) such that $Ag = -f$} \\[2pt]
374 & 343 & \sichB{} & Text vor Proposition 4.1.7 & Die Beschreibung $\Image{A} = \{f \in L^{2}(\R_{+}) \colon \int_{0}^{\infty} f(x)\dx \text{ exists}\}$ ist falsch: nach Remark a4-1.17 ist $f \in \Image{A}$ gleichwertig mit der Existenz von $\int_0^\infty T(t)f\dt$ in $L^2$, d.h. $x \mapsto \int_x^\infty f(y)\dy$ muss existieren und in $L^2(\R_+)$ liegen. Gegenbeispiel: $f(x) = \frac12(1+x)^{-3/2}$, $\int_0^\infty f = 1$, aber $\int_x^\infty f = (1+x)^{-1/2} \notin L^2$. Die Totalität von $(\Image{A})_+$ (positive Funktionen mit kompaktem Träger) bleibt richtig. & $\Image{A} = \{f \in L^{2}(\R_{+}) \colon x \mapsto \int_{x}^{\infty} f(y)\dy \text{ exists and belongs to } L^{2}(\R_{+})\}$ \\[2pt]
376 & 446 & \sichB{} & Beweis von Proposition 4.1.11 & $h \coloneq \int_{0}^{\infty} T(s)f_{0}\ds \in \Image{A}$ ist nicht begründet und nicht gemeint: nach Part-I, Remark a4-1.17 gilt $h \in D(A)$ mit $Ah=-f_0$ (so auch in der Aussage, Z. 435). & $\in D(A)$ \\[2pt]
384 & 847 & \sichB{} & Beweis von Corollary 4.2.7 & „Since $T(\tau)$ is a contraction on a reflexive Banach space“: Theorem 4.2.6 und das Korollar lassen $p = 1$ zu, $L^1(\mu)$ ist nicht reflexiv. Die Zerlegung $L^p = \Kern{(\Id-T(\tau))} \oplus \overline{\Image{(\Id-T(\tau))}}$ gilt trotzdem, da $T(\tau)$ das schwach kompakte Ordnungsintervall $[-e,e]$ invariant lässt und $\{g \colon \vert g\vert \le e\}$ total ist (Mean-Ergodic-Theorem). & Begründung ergänzen, z. B. „Since $T(\tau)$ is a contraction leaving the weakly compact order interval $[-e,e]$ invariant, $T(\tau)$ is mean ergodic and \dots{}“ – bitte prüfen \\[2pt]
386 & 1022 & \sichC{} & Beweis von Corollary 4.2.11 & „It follows that $T_j(t_0)$ is an irreducible kernel operator“: aus der Irreduzibilität der Halbgruppe folgt nicht automatisch die Irreduzibilität des Einzeloperators $T_j(t_0)$ (vgl. Rotationsgruppe: $T(\tau)=\Id$); Eq. (429) setzt aber genau diese voraus. Vermutlich ist das Argument über $P\sigma(A_j)\cap\im\R=\{0\}$ gemeint, es fehlt jedoch. & bitte prüfen (Begründung bzw. Zitat ergänzen) \\[2pt]
390 & 1202 & \sichA{} & Example 4.2.15 & Ergänzt korrekturen-mathe. Die Randbedingung $f(\frac{1}{2}) = 0$ ist mathematisch unvereinbar mit dem Rest: die angegebene Halbgruppe (Z. 1223–1226) und Resolvente (Z. 1229, $R(\lambda,A)f(\frac14)=0$) haben Generator mit $f(\frac14)=0$; die partielle Integration in Z. 1248 benötigt $f(\frac14)=0$; und die Eigenfunktion $h_0$ (Z. 1260/1261) erfüllt $h_0(\frac14)=0$, aber $h_0(\frac12)=1$, läge mit $f(\frac12)=0$ also nicht in $D(A)$. & $D(A) \coloneq \{f \in AC \colon f(\frac{1}{4}) = 0\}$ \\[2pt]
394 & 1431 & \sichC{} & Beweis von Theorem 4.3.1 & „Let $t \in \R_+$ and $s \in [-t,0]$“ passt nicht zur folgenden Fallunterscheidung: für $s \in [-t,0]$ ist $t+s \leq 0$ nur bei $s=-t$, der Fall $t+s\le 0$ wird aber auf $[-1,-t]$ behandelt (Z. 1433–1435). & $s \in [-1,0]$ \\[2pt]
394 & 1473 & \sichB{} & Text vor Proposition 4.3.3 & $R(\lambda,A) = R(1,S_\lambda)R(\lambda,A_{0})$ „for $\lambda > \omega$“: $\omega$ ist nirgends definiert; außerdem folgt $R(1,S_\lambda) \geq 0$ nicht schon aus der Positivität von $S_\lambda$, sondern erst aus der Neumann-Reihe für $r(S_\lambda) < 1$, also z. B. für $\lambda > \|\Phi\|$ (Z. 1369–1372). & {for $\lambda > \|\Phi\|$} und „Since $S_\lambda \geq 0$ and $\|S_\lambda\| < 1$, we have $R(1,S_\lambda) = \sum_n S_\lambda^n \geq 0$“ \\[2pt]
395 & 1502 & \sichC{} & Beweis von Lemma 4.3.4 & Corollary {\lbl{cor:c3-2.12}} liefert nur die Zyklizität von $\sigma_b$ (unter Polvoraussetzung); dass $\sigma_b(\Phi_\lambda) = \{s(\Phi_\lambda)\}$ ist (dominanter Spektralwert, beschränkter Generator $\Phi_\lambda$), steht in Corollary {\lbl{cor:c3-2.13}}. Zudem wird die Poleigenschaft erst im folgenden Satz genannt. & {(see Corollary~\lbl{cor:c3-2.13})}, Poleigenschaft vorziehen – bitte prüfen \\[2pt]
395 & 1515 & \sichA{} & Lemma 4.3.5 (i), Eq. (4.3.7) & Vorzeichenfehler: $u(t) = \int_{0}^{t} \eu^{\lambda(t-s)}g(s)\ds$ erfüllt $\lambda u - u' = -g$, nicht $= g$; für $t \in [-1,0]$, $\lambda>0$, $g\ge 0$ wäre $R(\lambda,A_0)g \le 0$, im Widerspruch zur Positivität. Laplace-Transformation von (4.3.1) ergibt $\int_{t}^{0} \eu^{\lambda(t-s)}g(s)\ds$. Ebenso Z. 1533 und 1538. & $(R(\lambda,A_{0})g)(t) = \int_{t}^{0} \eu^{\lambda(t-s)}g(s)\ds$ (ebenso Z. 1533, 1538) \\[2pt]
398 & 1696 & \sichC{} & Example 4.3.9, Lemma (ii) & Aus $\sigma(\Phi_{\lambda}) = \psi(\epsilon_{\lambda})\cdot\sigma(B)$ folgt „strictly decreasing“ nur, wenn $\psi \neq 0$ und $s(B) > 0$; für quasinilpotentes $B$ ist $s(\Phi_\mu) \equiv 0$. (Die Aussage (4.3.10) bleibt dann trivial richtig, da $\sigma(A)=\emptyset$.) & Voraussetzung $0 \neq \psi$, $s(B) > 0$ ergänzen – bitte prüfen \\[2pt]
398 & 1732 & \sichB{} & Example 4.3.10 & „unique solution $f \in L^{1}([-1,\infty),F)$“: Corollary 4.3.2 liefert nur eine Lösung in $L^{1}_{\mathrm{loc}}([-1,\infty),F)$; im Allgemeinen wachsen die Lösungen exponentiell (z. B. $s(A)>0$) und sind nicht integrierbar auf $[-1,\infty)$. & $f \in L^{1}_{\mathrm{loc}}([-1,\infty),F)$ \\[2pt]
398 & 1751 & \sichC{} & Beweis zu Example 4.3.10 & Die Rechnung wird nur für $0 < \mu < \lambda$ geführt, gebraucht wird strikte Monotonie auf ganz $\R$; die Abschätzung gilt wörtlich für alle $\mu < \lambda$ ($\eu^{(\mu-\lambda)s} > 1$ für $s<0$). Ebenso Z. 1820 ($\lambda > \mu > 0$). Zudem ist „Assume $s(\Phi_{\lambda}) > 0$ for all $\lambda$“ (Z. 1747) eine unbegründete Zusatzannahme. & {For $0 < x \in F$ and $\mu < \lambda$}; ebenso Z. 1820 $\lambda > \mu$; Z. 1747 begründen – bitte prüfen \\[2pt]
399 & 1783 & \sichB{} & Example 4.3.11 & Mit $F = L^{p}([0,1])$, $1<p<\infty$, ist $\Phi$ nicht beschränkt auf $E = L^1([-1,0],F)$ (Z. 1797 „easy to show“): für $k \equiv 1$, $\psi(s) = g(s)\1$ mit $g \in L^1 \setminus L^p$ ist $\Phi\psi = g\circ q \notin L^p$. Außerdem hat $L^p$ ($1<p<\infty$) nicht die Dunford-Pettis-Eigenschaft (Z. 1818). Alles stimmt für $p = 1$. & $F = L^{1}([0,1])$ – bitte am Original prüfen \\[2pt]
399 & 1818 & \sichB{} & Beweis zu Example 4.3.11 & „Since $E$ has the Dunford-Pettis property“: $\Phi_\lambda$ wirkt auf $F$; nötig ist die DP-Eigenschaft von $F$ (für $F=L^1$ richtig, für $L^p$, $1<p<\infty$, falsch; vgl. Z. 1783). & {Since $F$ has the Dunford-Pettis property} \\[2pt]
399 & 1858 & \sichC{} & Beweis zu Example 4.3.11 & Aus $r(\Phi_{\mu}) \geq \eu^{(\mu-\lambda)\epsilon}r(\Phi_{\lambda})$ folgt $r(\Phi_\mu) > r(\Phi_\lambda)$ nur bei $r(\Phi_{\lambda}) > 0$, was nicht gezeigt wird (folgt z. B. aus strikter Positivität des Kerns). Die strikte Ungleichung $\|(\Phi_{\mu})^n\| > \ldots$ in Z. 1846 ist nicht begründet; $\geq$ genügt. & $r(\Phi_\lambda) > 0$ begründen; Z. 1846 $\geq$ \\[2pt]
401 & 1994 & \sichA{} & Example 4.3.12 & „whose unique real solution $\lambda$ is $s(A)$“ ist falsch: $\det(\Id-\Phi_\lambda)=0$ kann mehrere reelle Lösungen haben. Gegenbeispiel $n=2$, $\beta_{11}=\beta_{22}=2$, $\beta_{12}=\beta_{21}=1$: $\Phi_\lambda = c(\lambda)\bigl(\begin{smallmatrix}2&1 1&2\end{smallmatrix}\bigr)$ mit Eigenwerten $3c(\lambda)$, $c(\lambda)$; beide Gleichungen $3c(\lambda)=1$, $c(\lambda)=1$ haben reelle Lösungen. Eindeutig ist nur die reelle Lösung von $s(\Phi_\lambda)=1$ (Cor. 4.3.8). & {whose largest real solution $\lambda$ is $s(A)$} \\[2pt]
402 & 2130 & \sichB{} & Notes, Section 2 & „false in $L^{\infty}$-spaces since \dots{} the identity operator satisfies Doeblin's condition“: Die Identität auf einem unendlichdimensionalen $L^\infty$ erfüllt weder (4.2.4) noch die Bedingung aus Part-II, Prop. b4-2.13 (disjunkte Indikatorfunktionen; dort ist Doeblin sogar hinreichend für Quasikompaktheit auf $C(K) \cong L^\infty$). In $L^1$ erfüllt die Identität (4.2.4) mit $\phi = \1$, $b=0$. & {false in $L^{1}$-spaces} \\
\end{longtable}
\end{document}
